\chapter{Мультиплексоры и демультиплексоры}
\label{ch:mux}

В главе~\ref{ch:coders} мы научились по номеру зажигать \emph{одну} линию из
$2^n$: дешифратор отвечает на вопрос «какой по счёту?». Но линия сама по себе
ничего не несёт. Остался вопрос: как по номеру не просто показать, а
\emph{пропустить} нужный сигнал, то есть выбрать один из нескольких источников
и отправить его дальше? Ответ даёт мультиплексор, и мы увидим, что внутри он
состоит из уже знакомого дешифратора.

\begin{buildbox}
\textbf{Задание.} Выберите, чьё значение показать на светодиоде: переключатель
\code{a} или переключатель \code{b}, по третьему переключателю \code{s}.

\textbf{В Logra.} Вкладка Обучение, книга «из блоков», глава «Выбор и
распределение», урок «Мультиплексор 2→1» (Конструктор). Нужен один блок
«Мультиплексор 2:1» из группы выбора.

\textbf{Ожидаемый результат.} При $s=0$ светодиод повторяет \code{a}, при
$s=1$ повторяет \code{b}. Остальное в этой главе --- объяснение того, что
внутри этого блока и во что его можно превратить.
\end{buildbox}

\section{Цифровой переключатель}

Представьте старый галетный переключатель: ручка поворачивается и соединяет
общий контакт с одним из нескольких. \term{Мультиплексор} (MUX) делает то же
самое с цифровыми сигналами: у него несколько информационных входов, один
выход и \term{адресные} (управляющие) входы, которые выбирают, какой из
информационных входов соединить с выходом.

\begin{figure}[H]
\centering
\begin{tikzpicture}
  % механический переключатель
  \foreach \i in {0,...,3} {
    \fill (0,1.5-\i) circle (2pt);
    \draw[wire] (-1,1.5-\i) node[left] {$x_\i$} -- (0,1.5-\i);
  }
  \fill (2,0) circle (2.5pt);
  \draw[very thick, FpgaOrange] (2,0) -- (0.08,0.5);
  \draw[wire] (2,0) -- (3,0) node[right] {$y$};
  \node[lbl] at (1.2,-2.3) {механический переключатель};
  % мультиплексор
  \begin{scope}[xshift=7.5cm]
    \draw[thick, fill=orange!15] (0,2) -- (1.6,1.2) -- (1.6,-1.2) -- (0,-2) -- cycle;
    \node[font=\sffamily\bfseries] at (0.8,0) {MUX};
    \foreach \i in {0,...,3} {
      \draw[wire] (-1,1.5-\i) node[left] {$x_\i$} -- (0,1.5-\i);
      \node[font=\sffamily\footnotesize, anchor=west] at (0.05,1.5-\i) {\i};
    }
    \draw[wire] (1.6,0) -- (2.6,0) node[right] {$y$};
    \draw[wire, FpgaBlue] (0.6,-1.7) -- (0.6,-2.6) node[below] {$s_1$};
    \draw[wire, FpgaBlue] (1.1,-1.45) -- (1.1,-2.6) node[below] {$s_0$};
    \node[lbl] at (0.8,-3.3) {цифровой мультиплексор 4:1};
  \end{scope}
  \draw[-{Stealth[length=4mm]}, line width=2pt, FpgaOrange] (4.2,0) -- (5.6,0);
\end{tikzpicture}
\caption{Мультиплексор --- это управляемый переключатель: при $s_1s_0 = 01$
выход соединён со входом $x_1$}
\end{figure}

Мультиплексор с $n$ адресными входами может выбирать из $2^n$
информационных входов. Его обозначают дробью: \textbf{2:1}, \textbf{4:1},
\textbf{8:1} (сколько входов : сколько выходов).

\begin{diybox}[Переключатель на четыре входа]
\textbf{Задание.} \stars{1} Повторите: в Конструкторе поставьте блок
«Мультиплексор 4:1» (параметр «Число входов» равен 4), к входам $D_0\ldots D_3$
подключите четыре разных источника (переключатели или константы), а вход $S$
сделайте двухбитным. \stars{2} Измените: сделайте так, чтобы на выход шёл
вход номер 2, и запишите, какой код нужен на $S$.

\textbf{В Logra.} Вкладки Конструктор и Диаграммы; образец --- пример
\code{mux-four}.

\textbf{Ожидаемый результат.} При $S=00,01,10,11$ выход равен $D_0$, $D_1$,
$D_2$, $D_3$; номер входа равен двоичному числу на $S$.

\textbf{Как проверить.} На вкладке Диаграммы переберите четыре значения $S$ и
сверьте с рис.~\ref{fig:mux41} и временной диаграммой ниже по тексту.
\end{diybox}

\section{Мультиплексор 2:1}

Самый простой мультиплексор имеет два входа $a$, $b$ и один адресный вход $s$:
если $s = 0$, на выход проходит $a$; если $s = 1$ --- проходит $b$.

\begin{figure}[H]
\centering
\begin{minipage}{0.24\textwidth}
\centering
\begin{tikzpicture}
  \draw[thick, fill=orange!15] (0,1.1) -- (1.1,0.6) -- (1.1,-0.6) -- (0,-1.1) -- cycle;
  \node[font=\sffamily\small] at (0.55,0) {MUX};
  \draw[wire] (-0.7,0.6) node[left] {$a$} -- (0,0.6); \node[font=\sffamily\footnotesize, anchor=west] at (0.02,0.6) {0};
  \draw[wire] (-0.7,-0.6) node[left] {$b$} -- (0,-0.6); \node[font=\sffamily\footnotesize, anchor=west] at (0.02,-0.6) {1};
  \draw[wire] (1.1,0) -- (1.7,0) node[right] {$y$};
  \draw[wire, FpgaBlue] (0.55,-0.85) -- (0.55,-1.5) node[below] {$s$};
\end{tikzpicture}
\end{minipage}
\begin{minipage}{0.21\textwidth}
\centering
\begin{tabular}{c|c}
\toprule
$s$ & $y$ \\
\midrule
0 & $a$ \\
1 & $b$ \\
\bottomrule
\end{tabular}

\medskip
$y = \overline{s}\,a + s\,b$
\end{minipage}
\begin{minipage}{0.52\textwidth}
\centering
\begin{tikzpicture}
  \node[gAND] (g1) at (3.7,1.2) {};
  \node[gAND] (g2) at (3.7,-1.2) {};
  \node[gOR] (o) at (6.0,0) {};
  \node[gNOT] (n) at (1.5,0) {};
  \draw[wire] (g1.input 1) -- (g1.input 1 -| 0,0) node[left] {$a$};
  \draw[wire] (g2.input 2) -- (g2.input 2 -| 0,0) node[left] {$b$};
  \draw[wire, FpgaBlue] (0,0) node[left, text=FpgaDark] {$s$} -- (0.5,0) coordinate (sj);
  \fill[FpgaBlue] (sj) circle (1.3pt);
  \draw[wire, FpgaBlue] (sj) -- (n.input);
  \draw[wire, FpgaBlue] (sj) |- (g2.input 1);
  \draw[wire] (n.output) -| ([xshift=-4mm]g1.input 2) -- (g1.input 2);
  \draw[wire] (g1.output) -| ([xshift=-4mm]o.input 1) -- (o.input 1);
  \draw[wire] (g2.output) -| ([xshift=-4mm]o.input 2) -- (o.input 2);
  \draw[wire] (o.output) -- ++(0.5,0) node[right] {$y$};
\end{tikzpicture}
\end{minipage}
\caption{Мультиплексор 2:1: обозначение, таблица, формула и схема на элементах}
\label{fig:mux21}
\end{figure}

Как работает схема: при $s = 0$ верхний элемент И «пропускает» $a$, а нижний
«закрыт» (на его входе 0). При $s = 1$ --- наоборот. Элемент ИЛИ собирает
результат.

\begin{diybox}[Мультиплексор из вентилей]
\textbf{Задание.} \stars{1} Соберите мультиплексор 2:1 из элементов НЕ, И, ИЛИ
по рис.~\ref{fig:mux21}. \stars{2} Замените все три вида элементов элементами
И-НЕ (глава~\ref{ch:gates}) и сравните число элементов.

\textbf{В Logra.} Вкладка Обучение, глава «Выбор и распределение», урок
«Мультиплексор из НЕ, И, ИЛИ»; для сравнения --- готовый блок и пример
\code{mux}.

\textbf{Ожидаемый результат.} Ваша схема и готовый блок дают одинаковый выход
при всех восьми комбинациях $a,b,s$.

\textbf{Как проверить.} Диаграммы: подайте на обе схемы одни и те же входы и
сравните выходы; RTL: сравните свой рисунок с рисунком блока.
\end{diybox}

\section{Мультиплексор 4:1}

У мультиплексора 4:1 два адресных входа $s_1 s_0$ задают номер входа от 0 до 3.
\[
  y = \overline{s_1}\,\overline{s_0}\,x_0 + \overline{s_1}\,s_0\,x_1 + s_1\,\overline{s_0}\,x_2 + s_1\,s_0\,x_3.
\]
Присмотритесь к формуле: множители при $x_i$ --- это ровно выходы
\emph{дешифратора} 2$\to$4 из главы~\ref{ch:coders}! Значит, мультиплексор
состоит из дешифратора адреса, четырёх элементов И и одного ИЛИ.

\begin{figure}[H]
\centering
\begin{tikzpicture}
  \draw[thick, fill=FpgaLight] (0,-5.5) rectangle (1.6,0.8);
  \node[font=\sffamily\bfseries] at (0.8,-2.35) {DC};
  \draw[wire, FpgaBlue] (-0.8,-1.5) node[left, text=FpgaDark] {$s_1$} -- (0,-1.5);
  \draw[wire, FpgaBlue] (-0.8,-3.2) node[left, text=FpgaDark] {$s_0$} -- (0,-3.2);
  \foreach \i in {0,...,3} {
    \node[gAND] (g\i) at (5.0,-\i*1.5-0.3) {};
    \draw[wire] (1.6,0 |- g\i.input 2) -- (g\i.input 2);
    \node[font=\sffamily\footnotesize, anchor=south east, inner sep=1pt] at (1.6,0 |- g\i.input 2) {\i};
    \draw[wire] (g\i.input 1) -- ++(-1.8,0) node[left, font=\small] {$x_\i$};
  }
  \node[gOR, logic gate inputs=nnnn, minimum height=18mm] (o) at (8.9,-2.55) {}; % fig-ok: 4-входовой вентиль 18 мм
  \draw[wire] (g0.output) -| ([xshift=-5mm]o.input 1) -- (o.input 1);
  \draw[wire] (g1.output) -| ([xshift=-10mm]o.input 2) -- (o.input 2);
  \draw[wire] (g2.output) -| ([xshift=-10mm]o.input 3) -- (o.input 3);
  \draw[wire] (g3.output) -| ([xshift=-5mm]o.input 4) -- (o.input 4);
  \draw[wire] (o.output) -- ++(0.7,0) node[right] {$y$};
\end{tikzpicture}
\caption{Мультиплексор 4:1 = дешифратор адреса + элементы И + элемент ИЛИ}
\label{fig:mux41}
\end{figure}

\begin{figure}[H]
\centering
\begin{tikztimingtable}[timing/slope=0.05, xscale=0.9, yscale=1.4, timing/font=\sffamily\small, timing/rowdist=3.2ex]
  $x_0$     & 16{1L 1H} \\
  $x_1$     & 8{2L 2H} \\
  $x_2$     & 32H \\
  $x_3$     & 32L \\
  $s_1s_0$  & 8D{00} 8D{01} 8D{10} 8D{11} \\
  $y$       & 4{1L 1H} 2{2L 2H} 8H 8L \\
\end{tikztimingtable}
\caption{Мультиплексор 4:1 по очереди передаёт на выход разные входные
сигналы}
\end{figure}

\begin{diybox}[Мультиплексор 4:1]
\textbf{Задание.} \stars{1} Пройдите урок и соберите мультиплексор 4:1
блоком. \stars{2} Соберите его ещё раз «по рис.~\ref{fig:mux41}»: дешифратор
адреса, четыре элемента И, один элемент ИЛИ.

\textbf{В Logra.} Вкладка Обучение, глава «Выбор и распределение», урок
«Мультиплексор 4→1»; для второго задания --- урок «Дешифратор» как источник
дешифратора.

\textbf{Ожидаемый результат.} Обе схемы дают одну и ту же таблицу: выход равен
входу с номером, заданным на $s_1s_0$.

\textbf{Как проверить.} Диаграммы: переберите все $4$ адреса при разных
значениях $x_0\ldots x_3$ (достаточно по два набора); RTL: блок и сборка
должны отличаться только видом, не логикой.
\end{diybox}

\section{Дерево мультиплексоров}

Большой мультиплексор можно собрать из маленьких. Мультиплексор 4:1 --- это
три мультиплексора 2:1: первый уровень выбирает по младшему биту адреса,
второй --- по старшему. Мультиплексор 8:1 --- дерево из семи мультиплексоров
2:1 в три уровня, и так далее.

\newcommand{\muxsym}[3]{% имя, x, y
  \draw[thick, fill=orange!15] (#2,#3+0.55) -- (#2+0.7,#3+0.3) -- (#2+0.7,#3-0.3) -- (#2,#3-0.55) -- cycle;
  \coordinate (#1-i0) at (#2,#3+0.3); \coordinate (#1-i1) at (#2,#3-0.3);
  \coordinate (#1-o) at (#2+0.7,#3); \coordinate (#1-s) at (#2+0.35,#3-0.42);}

\begin{figure}[H]
\centering
\begin{tikzpicture}
  \muxsym{m0}{0}{1}
  \muxsym{m1}{0}{-1}
  \muxsym{m2}{2}{0}
  \foreach \m/\a/\b in {m0/x_0/x_1, m1/x_2/x_3} {
    \draw[wire] (\m-i0) -- ++(-0.6,0) node[left] {$\a$};
    \draw[wire] (\m-i1) -- ++(-0.6,0) node[left] {$\b$};
  }
  \draw[wire] (m0-o) -- ++(0.4,0) |- (m2-i0);
  \draw[wire] (m1-o) -- ++(0.4,0) |- (m2-i1);
  \draw[wire] (m2-o) -- ++(0.6,0) node[right] {$y$};
  \draw[wire, FpgaBlue] (m0-s) -- (0.35,-0.575);
  \fill[FpgaBlue] (0.35,0) circle (1.3pt);
  \draw[wire, FpgaBlue] (0.35,0) -- (-0.9,0) node[left, text=FpgaDark] {$s_0$};
  \draw[wire, FpgaBlue] (m2-s) -- (2.35,-1.2) node[below, text=FpgaDark] {$s_1$};
  \node[lbl] at (0.35,-2.4) {уровень 1};
  \node[lbl] at (2.35,-2.4) {уровень 2};
\end{tikzpicture}
\caption{Мультиплексор 4:1 из трёх мультиплексоров 2:1}
\label{fig:muxtree}
\end{figure}

\begin{diybox}[Дерево мультиплексоров]
\textbf{Задание.} \stars{1} Соберите мультиплексор 4:1 из трёх блоков
«Мультиплексор 2:1» по рис.~\ref{fig:muxtree}. \stars{2} Добавьте ещё четыре
блока и получите мультиплексор 8:1; запишите, какие разряды адреса подаются на
какой уровень.

\textbf{В Logra.} Вкладки Конструктор и RTL; образец --- пример \code{mux-four}
(сравнить с готовым блоком).

\textbf{Ожидаемый результат.} Дерево и готовый блок дают одинаковый выход при
всех адресах; младший разряд адреса $s_0$ идёт на первый уровень.

\textbf{Как проверить.} Диаграммы: подайте на $x_0\ldots x_3$ постоянные
значения $0,1,0,1$ и переберите адрес --- выход должен чередоваться
$0,1,0,1$.
\end{diybox}

\section{Мультиплексор как универсальный логический элемент}

Вот удивительное свойство мультиплексора. Подадим на адресные входы
\emph{переменные} $a$ и $b$, а на информационные входы --- \emph{константы}
0 и 1, взятые из правого столбца таблицы истинности какой-либо функции. Тогда
мультиплексор будет выбирать нужную строку таблицы и выдавать значение
функции! Любую функцию от $n$ переменных можно реализовать одним
мультиплексором $2^n$:1.

\begin{figure}[H]
\centering
\begin{tikzpicture}
  \draw[thick, fill=orange!15] (0,2) -- (1.6,1.2) -- (1.6,-1.2) -- (0,-2) -- cycle;
  \node[font=\sffamily\bfseries] at (0.8,0) {MUX};
  \foreach \i/\v/\ab in {0/0/00, 1/1/01, 2/1/10, 3/0/11} {
    \draw[wire] (-0.8,1.5-\i) -- (0,1.5-\i);
    \node[draw=FpgaGreen, fill=green!8, thick, minimum size=6mm, font=\ttfamily] at (-1.1,1.5-\i) {\v};
    \node[font=\sffamily\footnotesize, anchor=west] at (0.05,1.5-\i) {\i};
    \node[lbl, anchor=east] at (-1.5,1.5-\i) {$ab=\ab$};
  }
  \draw[wire] (1.6,0) -- (2.6,0) node[right] {$y = a\oplus b$};
  \draw[wire, FpgaBlue] (0.6,-1.7) -- (0.6,-2.6) node[below, text=FpgaDark] {$a$};
  \draw[wire, FpgaBlue] (1.1,-1.45) -- (1.1,-2.6) node[below, text=FpgaDark] {$b$};
  \node[lbl, align=center, text=FpgaGreen] at (-1.1,2.3) {таблица\\истинности};
  \node[font=\sffamily\small, align=left, anchor=west] at (5.0,0.8) {Подставьте \code{0001} --- получится И,};
  \node[font=\sffamily\small, align=left, anchor=west] at (5.0,0.2) {\code{0111} --- ИЛИ, \code{1110} --- И-НЕ\ldots};
  \node[font=\sffamily\small, align=left, anchor=west] at (5.0,-0.4) {Схема та же --- меняются только};
  \node[font=\sffamily\small, align=left, anchor=west] at (5.0,-1.0) {четыре бита слева.};
\end{tikzpicture}
\caption{Мультиплексор 4:1 с константами на входах вычисляет любую функцию
двух переменных (здесь --- Исключающее ИЛИ)}
\label{fig:mux_lut}
\end{figure}

\begin{ideabox}[Запомните этот рисунок]
Четыре константы слева --- это просто столбец значений из таблицы истинности.
Меняя их, мы меняем функцию, а схема остаётся прежней. Если эти биты хранить в
маленькой памяти, получится «программируемая» логика.

\lookahead{такой мультиплексор с памятью называется LUT; из тысяч таких блоков
состоит ПЛИС (глава~\ref{ch:inside}).}
\end{ideabox}

\begin{diybox}[Любая функция на одном мультиплексоре]
\textbf{Задание.} \stars{1} Реализуйте $A\oplus B$ на мультиплексоре 4:1 по
рис.~\ref{fig:mux_lut}. \stars{2} Реализуйте $F=A\oplus B\oplus C$ (задача~5
этой главы): на адрес $A,B$, а на входы данных --- константы и $C$ или
$\overline{C}$. \stars{3} Реализуйте на одном блоке «2 из 3»
(мажоритарную функцию из главы~\ref{ch:comb}).

\textbf{В Logra.} Вкладки Конструктор, Диаграммы, RTL; образец --- пример
\code{mux-four}.

\textbf{Ожидаемый результат.} Таблица выхода совпадает с таблицей истинности
функции при всех комбинациях; меняя только константы на входах данных, вы
получаете И, ИЛИ, И-НЕ без перестройки схемы.

\textbf{Как проверить.} Диаграммы: переберите все комбинации и сравните с
таблицей; RTL: убедитесь, что структура одна и та же для разных функций.
\end{diybox}

\section{Демультиплексор}

\term{Демультиплексор} (DEMUX) выполняет обратное действие: у него один
информационный вход $x$ и $2^n$ выходов. Адрес $s$ определяет, на какой из
выходов будет передан входной сигнал; на остальных выходах --- 0.

\begin{figure}[H]
\centering
\begin{minipage}{0.35\textwidth}
\centering
\begin{tikzpicture}
  \draw[thick, fill=orange!15] (0,1.2) -- (1.6,2) -- (1.6,-2) -- (0,-1.2) -- cycle;
  \node[font=\sffamily\small\bfseries] at (0.8,0) {DEMUX};
  \draw[wire] (-0.8,0) node[left] {$x$} -- (0,0);
  \foreach \i in {0,...,3} {
    \draw[wire] (1.6,1.5-\i) -- (2.4,1.5-\i) node[right] {$y_\i$};
    \node[font=\sffamily\footnotesize, anchor=east] at (1.55,1.5-\i) {\i};
  }
  \draw[wire, FpgaBlue] (0.5,-1.45) -- (0.5,-2.6) node[below, text=FpgaDark] {$s_1$};
  \draw[wire, FpgaBlue] (1.0,-1.7) -- (1.0,-2.6) node[below, text=FpgaDark] {$s_0$};
\end{tikzpicture}
\end{minipage}\hfill
\begin{minipage}{0.6\textwidth}
\centering
\begin{tabular}{cc|cccc}
\toprule
$s_1$ & $s_0$ & $y_0$ & $y_1$ & $y_2$ & $y_3$ \\
\midrule
0 & 0 & $x$ & 0 & 0 & 0 \\
0 & 1 & 0 & $x$ & 0 & 0 \\
1 & 0 & 0 & 0 & $x$ & 0 \\
1 & 1 & 0 & 0 & 0 & $x$ \\
\bottomrule
\end{tabular}

\medskip
$y_0 = x\,\overline{s_1}\,\overline{s_0}, \ \ y_1 = x\,\overline{s_1}\,s_0, \ \ \ldots$
\end{minipage}
\caption{Демультиплексор 1:4}
\end{figure}

Сравните формулы с дешифратором со входом разрешения
(рис.~\ref{fig:dc24}): они совпадают, если вход $E$ назвать $x$. Поэтому
\textbf{демультиплексор --- это дешифратор, у которого на вход разрешения
подаётся информационный сигнал}. В справочниках эти две микросхемы часто
объединяют под одним названием.

\begin{diybox}[Демультиплексор]
\textbf{Задание.} \stars{1} Соберите демультиплексор: одна кнопка, четыре
лампы, номер лампы задают два переключателя. \stars{2} Постройте то же самое
из дешифратора со входом разрешения $E$, подав кнопку на $E$.

\textbf{В Logra.} Вкладка Обучение, глава «Выбор и распределение», урок
«Демультиплексор»; пример \code{demux-four}; для второго задания --- урок
«Дешифратор».

\textbf{Ожидаемый результат.} Нажатая кнопка зажигает только одну лампу,
выбранную адресом; остальные остаются выключенными.

\textbf{Как проверить.} Диаграммы: при адресе 10 выход $y_2$ повторяет кнопку,
остальные равны 0; RTL: блок и дешифратор с $E$ дают одинаковые диаграммы.
\end{diybox}

\section{Передача по одной линии}

Вместе мультиплексор и демультиплексор позволяют передать несколько сигналов
по \emph{одному} проводу. Адрес быстро перебирает значения $0, 1, 2, 3, 0,
1\ldots$, и в каждый момент по линии идёт сигнал одного из источников. На
приёмной стороне демультиплексор с таким же адресом раскладывает их по своим
выходам. Это \term{временное разделение каналов}: так работают телефонные
линии, шины компьютеров, а в дронах --- приёмник пульта, передающий по одному
проводу сигналы всех каналов (протоколы PPM, SBUS).

\begin{figure}[H]
\centering
\begin{tikzpicture}
  \draw[thick, fill=orange!15] (0,1.3) -- (1.2,0.7) -- (1.2,-0.7) -- (0,-1.3) -- cycle;
  \node[font=\sffamily\footnotesize\bfseries] at (0.6,0) {MUX};
  \foreach \i/\t in {0/газ, 1/крен, 2/тангаж, 3/рыскание}
    \draw[wire] (-0.6,0.9-\i*0.6) node[left, font=\sffamily\small] {\t} -- (0,0.9-\i*0.6);
  \draw[very thick, FpgaOrange] (1.2,0) -- node[above, font=\sffamily\small, text=FpgaDark] {одна линия} (5.0,0);
  \begin{scope}[xshift=5cm]
  \draw[thick, fill=orange!15] (0,0.7) -- (1.2,1.3) -- (1.2,-1.3) -- (0,-0.7) -- cycle;
  \node[font=\sffamily\footnotesize\bfseries] at (0.6,0) {DEMUX};
  \foreach \i/\t in {0/газ, 1/крен, 2/тангаж, 3/рыскание}
    \draw[wire] (1.2,0.9-\i*0.6) -- (1.8,0.9-\i*0.6) node[right, font=\sffamily\small] {\t};
  \end{scope}
  \node[block, minimum width=20mm, minimum height=7mm] (ctr) at (3.1,-2.1) {Счётчик\\\footnotesize адреса};
  \draw[arr, FpgaBlue] (ctr.west) -| node[pos=0.25, below, lbl] {$s$} (0.6,-1.05);
  \draw[arr, FpgaBlue] (ctr.east) -| node[pos=0.25, below, lbl] {$s$} (5.6,-1.05);
\end{tikzpicture}
\caption{Четыре канала пульта по одному проводу}
\end{figure}

\begin{diybox}[Охранная панель: четыре зоны через один мультиплексор]
\textbf{Задание.} \stars{1} Соберите «мультиплексор 4:1 на четыре датчика»:
$x_0\ldots x_3$ --- датчики зон, адрес $s_1s_0$ задают два переключателя,
выход показывает состояние выбранной зоны. \stars{2} Добавьте демультиплексор
и покажите номер выбранной зоны отдельной лампой. \stars{3} Замените
переключатели адреса счётчиком, чтобы зоны опрашивались по очереди
(счётчики --- в главе~\ref{ch:regs}; пока можно менять адрес вручную).

\textbf{В Logra.} Вкладки Конструктор и Диаграммы; образцы --- урок
«Мультиплексор 4→1» и пример \code{mux-four}.

\textbf{Ожидаемый результат.} Выход мультиплексора равен датчику выбранной зоны
и больше ни от чего не зависит.

\textbf{Этап проекта «Охранная панель»: этап~4б.} Четыре зоны через один
мультиплексор. Критерий приёмки: при всех $4\times 2=8$ сочетаниях «адрес
зоны --- датчик этой зоны открыт/закрыт» выход равен датчику выбранной зоны,
а изменение датчиков других зон на выход не влияет (проверить на вкладке
Диаграммы).
\end{diybox}

\section{Многоразрядные мультиплексоры}

Часто нужно выбирать не между битами, а между \emph{числами}. Для этого
несколько мультиплексоров ставят параллельно с общим адресом. Например,
8-разрядный мультиплексор 2:1 --- это восемь мультиплексоров 2:1, у
которых соединены адресные входы. На схемах такой узел рисуют одним символом
с шинами.

\begin{figure}[H]
\centering
\begin{tikzpicture}
  \draw[thick, fill=orange!15] (0,1.1) -- (1.1,0.6) -- (1.1,-0.6) -- (0,-1.1) -- cycle;
  \node[font=\sffamily\small] at (0.55,0) {MUX};
  \draw[bus] (-1.4,0.6) node[left, text=FpgaDark] {$A$} -- (0,0.6);
  \draw[bus] (-1.4,-0.6) node[left, text=FpgaDark] {$B$} -- (0,-0.6);
  \draw[bus] (1.1,0) -- (2.4,0) node[right, text=FpgaDark] {$Y$};
  \foreach \x/\y in {-0.8/0.6, -0.8/-0.6, 1.7/0} {
    \draw[thick] (\x-0.1,\y-0.15) -- (\x+0.1,\y+0.15);
    \node[font=\sffamily\footnotesize, above] at (\x,\y+0.1) {8};
  }
  \draw[wire, FpgaBlue] (0.55,-0.85) -- (0.55,-1.5) node[below, text=FpgaDark] {$s$};
  \node[font=\sffamily\small, anchor=west] at (4.0,0.3) {$s=0$: \ $Y = A$};
  \node[font=\sffamily\small, anchor=west] at (4.0,-0.3) {$s=1$: \ $Y = B$};
\end{tikzpicture}
\caption{Мультиплексор 8-разрядных чисел. Косая черта с цифрой показывает число проводов в шине}
\end{figure}

\begin{ideabox}
\textbf{Мультиплексор} выбирает один из $2^n$ входов по адресу. Он состоит из
дешифратора адреса и элементов И-ИЛИ и может реализовать любую функцию, если
подать на входы константы. \textbf{Демультиплексор} передаёт один вход на
выбранный выход --- это дешифратор с информационным входом разрешения.
\end{ideabox}

\begin{diybox}[Выбор между числами]
\textbf{Задание.} \stars{1} Выберите одно из трёх 4-разрядных чисел $a,b,c$
по двухбитному коду $s$ (при $s=3$ пусть выбирается $c$). \stars{2} Опишите
выбор между двумя 8-разрядными шинами одной строкой в Коде:
\code{assign y = s ? b : a;}, где \code{a} и \code{b} --- по 8 бит.
\stars{3} Добавьте третий источник и сравните с блоком «Мультиплексор 4:1»
с параметром «Ширина данных», равным~8.

\textbf{В Logra.} Вкладка Обучение, глава «Выбор и распределение», урок
«Выбор из трёх источников»; вкладка Код для второго и третьего заданий; пример
\code{mux}.

\textbf{Ожидаемый результат.} Один блок (одна строка кода) переключает сразу
все разряды; адресный вход $s$ остаётся однобитным (или двухбитным).

\textbf{Как проверить.} Диаграммы: подайте два разных числа, например \code{0xA5}
и \code{0x3C}, и переключайте $s$ --- выход должен целиком повторять выбранное
число; RTL: один символ мультиплексора с шинами, как на рисунке выше.
\end{diybox}

\section{Примеры на Verilog}
\label{sec:mux-verilog}

На Verilog мультиплексор --- это одна строка или короткий \code{case}.
Пять примеров ниже --- законченные модули из папки \code{code/b5\_mux/}; у
каждого есть тестбенч \code{tb\_*.v}, который сравнивает варианты друг с другом
и с ожидаемым значением и печатает \code{OK}. Запуск в Logra: вкладка
\emph{Код} $\to$ \emph{Синтезировать} (\code{Ctrl+Enter}) $\to$ \emph{RTL}
$\to$ \emph{Диаграммы} (раздел~\ref{sec:signals-verilog}).

\subsection{Пример 1. Мультиплексор 2:1 тремя способами}

Одна функция $y=\overline{s}\,a+s\,b$ записана оператором \code{?:}, через
\code{case} и вентилями из рис.~\ref{fig:mux21}.

\vfile{\codepath/b5_mux/mux2.v}{code/b5\_mux/mux2.v}

\textbf{Как читается.}
\begin{itemize}
  \item \code{assign y = s ? b : a;} --- «если $s$, то $b$, иначе $a$». Самый короткий способ.
  \item \code{always @(*)} + \code{case (s)} --- то же через оператор выбора. Выход объявлен \code{reg}, но схема остаётся комбинационной. \code{default} гарантирует, что \code{y} получает значение при любом $s$.
  \item \code{mux2\_gates} --- схема с рис.~\ref{fig:mux21}: \code{wire ns = \textasciitilde s;} и два И, объединённые ИЛИ.
\end{itemize}
\textbf{Проверка.} \code{tb\_mux2.v} подключает все три модуля к одним входам и
перебирает восемь комбинаций $s,a,b$: выходы должны совпасть. В Logra
синтезируйте каждый модуль и сравните \emph{RTL}: синтезатор может получить из
трёх записей близкие схемы.

\textbf{Допиши сам.} Добавьте вход разрешения \code{en}: при \code{en} $=0$ выход
равен нулю.

\textbf{Образец.}
\begin{vcode}
module mux2_en (
    input  a, b, s, en,
    output y
);
    assign y = en ? (s ? b : a) : 1'b0;   // вложенный ?:
endmodule
\end{vcode}

\subsection{Пример 2. Мультиплексор 4:1: case, индекс и дерево}

Для четырёх входов есть три удобные записи. Дерево повторяет
рис.~\ref{fig:muxtree}: два мультиплексора первого уровня и один второго.

\vfile{\codepath/b5_mux/mux4.v}{code/b5\_mux/mux4.v}

\textbf{Как читается.}
\begin{itemize}
  \item \code{input [3:0] x} --- четыре информационных входа как одна шина, \code{input [1:0] s} --- адрес.
  \item \code{x[s]} --- выбор бита шины по номеру: сам по себе мультиплексор.
  \item \code{mux4\_tree}: \code{m0} выбирает в паре ($x_0,x_1$) по \emph{младшему} биту адреса, \code{m1} --- в паре ($x_2,x_3$), \code{y} выбирает между парами по старшему. Порядок битов адреса важен: поменяйте \code{s[0]} и \code{s[1]}, и входы 1 и 2 поменяются местами.
\end{itemize}
\textbf{Проверка.} \code{tb\_mux4.v} перебирает все 16 значений \code{x} и все
4 адреса и сверяет три модуля с \code{x[s]}. В Logra: подайте \code{x} $=$
\code{0101} и меняйте \code{s} --- выход должен чередоваться $1,0,1,0$.

\textbf{Допиши сам.} Постройте мультиплексор 8:1 (три адресных бита).

\textbf{Образец.}
\begin{vcode}
module mux8 (
    input  [7:0] x,
    input  [2:0] s,
    output       y
);
    assign y = x[s];     // индекс по трёхбитному адресу: 8 входов
endmodule
\end{vcode}
Попробуйте также собрать его деревом из двух мультиплексоров 4:1 и одного 2:1.

\subsection{Пример 3. Демультиплексор 1:4}

Информационный вход \code{x} должен попасть на выход с номером \code{s}, на
остальных --- ноль. Первая запись --- через \code{case}, вторая --- сдвигом.

\vfile{\codepath/b5_mux/demux4.v}{code/b5\_mux/demux4.v}

\textbf{Как читается.}
\begin{itemize}
  \item \code{y = 4'b0000;} в начале блока --- \emph{значение по умолчанию}: дальше выбранный разряд перезаписывается. Так \code{y} присваивается на любой ветке, и защёлки нет, даже если в \code{case} перечислены не все значения.
  \item \code{y[2] = x;} --- присваивание одному разряду шины.
  \item \code{\{3'b000, x\} << s} --- число \code{000x} сдвигается влево на $s$ позиций: если $x=1$, единица окажется в разряде $s$. Это дешифратор со входом разрешения $x$.
\end{itemize}
\textbf{Проверка.} \code{tb\_demux4.v} перебирает $x$ и все адреса и сверяет оба
модуля. В Logra: при \code{s} $=10$ выход \code{y[2]} повторяет \code{x},
остальные равны нулю.

\textbf{Допиши сам.} Добавьте вход разрешения \code{en}: при \code{en} $=0$ все
выходы равны нулю.

\textbf{Образец.}
\begin{vcode}
module demux4_en (
    input        x, en,
    input  [1:0] s,
    output [3:0] y
);
    assign y = en ? ({3'b000, x} << s) : 4'b0000;
endmodule
\end{vcode}

\subsection{Пример 4. Мультиплексор как универсальный вентиль}

XOR получается из мультиплексора 2:1, если на входы подать $b$ и $\overline{b}$,
а адресом сделать $a$. Любая функция двух переменных --- один мультиплексор 4:1
с четырьмя константами; в Verilog константы собирает параметр.

\vfile{\codepath/b5_mux/mux_logic.v}{code/b5\_mux/mux\_logic.v}

\textbf{Как читается.}
\begin{itemize}
  \item \code{a ? \textasciitilde b : b} --- при $a=0$ выход равен $b$, при $a=1$ равен $\overline b$: исключающее ИЛИ.
  \item \code{parameter [3:0] TT} --- столбец таблицы истинности как параметр модуля; \code{TT[\{a,b\}]} выбирает бит по двухбитному адресу. Один и тот же модуль с \code{TT = 4'b1000} даёт И, с \code{4'b1110} --- ИЛИ, с \code{4'b0110} --- исключающее ИЛИ. Это LUT из рис.~\ref{fig:mux_lut}.
  \item \code{mux\_lut2 \#(.TT(4'b1000)) u (\ldots)} --- подключение модуля с заданным параметром.
\end{itemize}
\textbf{Проверка.} \code{tb\_mux\_logic.v} проверяет XOR, И, ИЛИ, а затем создаёт
сразу 16 экземпляров с \code{TT} от 0 до 15 и убеждается, что для каждой строки
таблицы выход равен соответствующему биту параметра. В Logra синтезируйте
\code{xor\_on\_mux} и сравните \emph{RTL} с рис.~\ref{fig:mux_lut}.

\textbf{Допиши сам.} Реализуйте на одном мультиплексоре 4:1 функцию «2 из 3»
$F=ab+ac+bc$ (адрес $\{a,b\}$, на входах данных --- константы и $c$).

\textbf{Образец.}
\begin{vcode}
module maj3_mux (
    input  a, b, c,
    output y
);
    // d[0] при ab=00 -> 0;  d[1] при 01 -> c;  d[2] при 10 -> c;  d[3] при 11 -> 1
    wire [3:0] d = {1'b1, c, c, 1'b0};
    assign y = d[{a, b}];
endmodule
\end{vcode}

\subsection{Пример 5. Мультиплексоры 8-разрядных шин}

Выбор между числами записывается так же, как между битами: ширина шин ни на что
не влияет.

\vfile{\codepath/b5_mux/mux_bus8.v}{code/b5\_mux/mux\_bus8.v}

\textbf{Как читается.} \code{y = s ? b : a;} при 8-разрядных \code{a} и \code{b}
переключает сразу все восемь разрядов: синтезатор ставит восемь мультиплексоров
с общим адресом. Версия 4:1 --- тот же \code{case}, но по двухбитному адресу
и с шинами вместо битов; \code{default} закрывает $s=3$.

\textbf{Проверка.} \code{tb\_mux\_bus8.v} подаёт числа \code{0xA5}, \code{0x3C},
\code{0xFF}, \code{0x00} и проверяет, что выход целиком повторяет выбранное
число. В Logra подайте \code{a} $=$ \code{A5}, \code{b} $=$ \code{3C} и
переключайте \code{s}; на \emph{RTL} это один символ мультиплексора с шинами,
как на рисунке «Мультиплексор 8-разрядных чисел».

\textbf{Допиши сам.} Выберите одно из трёх 4-разрядных чисел $a,b,c$ по
двухбитному коду $s$ (при $s=3$ выбирается $c$) --- задание из «Сделай сам».

\textbf{Образец.}
\begin{vcode}
module mux3_bus4 (
    input  [3:0] a, b, c,
    input  [1:0] s,
    output reg [3:0] y
);
    always @(*) begin
        case (s)
            2'd0:    y = a;
            2'd1:    y = b;
            default: y = c;    // s = 2 и s = 3
        endcase
    end
endmodule
\end{vcode}

\section{Типичные ошибки}

\begin{warnbox}[Типичные ошибки]
\begin{itemize}
  \item \textbf{Путают порядок адресных разрядов.} В коде $s_1s_0=01$ старший разряд~--- $s_1$, поэтому выбирается $x_1$, а не $x_2$. Если в Конструкторе поменять провода $s_1$ и $s_0$, зоны 1 и 2 поменяются местами, а схема внешне «работает».
  \item \textbf{Считают, что адресных входов $2^n$.} Мультиплексор $2^n{:}1$ имеет $n$ адресных входов: у $4{:}1$ их два, а не четыре, у $8{:}1$ --- три.
  \item \textbf{Ждут, что демультиплексор «запомнит» значение.} У него нет памяти: на невыбранных выходах всегда 0, и если адрес меняется, предыдущий выход сразу гаснет.
  \item \textbf{Рассинхронизируют адреса на передатчике и приёмнике.} При передаче по одной линии мультиплексор и демультиплексор должны получать \emph{один и тот же} адрес в один и тот же момент; сдвиг хотя бы на один шаг раскладывает каналы по чужим выходам.
  \item \textbf{Verilog: \code{case} без \code{default} (защёлка).}
    \code{always @(*) case (s) 2'd0: y = a; 2'd1: y = b; endcase} не задаёт
    $y$ при $s=2,3$, и синтезатор вставляет защёлку, которая «помнит» прежнее
    значение. Либо добавьте \code{default}, либо задайте значение по умолчанию в
    начале блока (\code{y = 1'b0;}).
  \item \textbf{Verilog: неполный список чувствительности.}
    \code{always @(a or b)} без \code{s} в симуляции не реагирует на смену
    адреса, хотя синтезированная схема реагирует. Пишите \code{always @(*)}.
  \item \textbf{Verilog: \code{<=} в комбинационном блоке.} В
    \code{always @(*)} используйте блокирующее \code{=}, в тактируемом
    \code{always @(posedge clk)} --- неблокирующее \code{<=}. Смешивание даёт
    расхождение между симуляцией и платой.
\end{itemize}
\end{warnbox}

\section{Вопросы и задачи}
\label{sec:tasks-mux}

Ответы собраны в приложении~\ref{app:answers}.

\begin{taskbox}[Задачи к главе]
\begin{enumerate}
  \item \stars{1}~Сколько адресных входов у мультиплексора $8:1$? А у $16:1$?
  \item \stars{1}~Сколько мультиплексоров $2:1$ нужно, чтобы собрать мультиплексор $8:1$?
  \item \stars{1}~Сколько мультиплексоров $2:1$ нужно, чтобы выбирать между двумя 8-разрядными шинами?
  \item \stars{1}~У демультиплексора $1:4$ на адресе 10 данные равны 1. Какие значения имеют четыре выхода?
  \item \stars{2}~Реализуйте функцию $F = A\oplus B\oplus C$ на одном мультиплексоре $4:1$, подав на адресные входы $A$ и $B$. Что нужно подать на информационные входы?
  \item \stars{2}~Почему мультиплексор называют универсальным логическим элементом? Какие четыре константы на входах $D_0\ldots D_3$ (при адресе $AB$) дают функции И и ИЛИ?
  \item \stars{2}~\emph{Найди ошибку.} Студент реализует $A\oplus B$ на мультиплексоре $4:1$ с адресом $AB$ и подаёт на $D_0,D_1,D_2,D_3$ константы $0,1,1,1$. Какая функция получилась на самом деле и какие константы нужны?
  \item \stars{2}~\emph{Найди ошибку.} В системе передачи по одной линии (раздел «Передача по одной линии») счётчик адреса приёмника запущен на один такт позже счётчика передатчика. Что окажется на выходе «газ», «крен», «тангаж», «рыскание»?
  \item \stars{3}~\textbf{Проект: этап 4б «Охранная панель».} Четыре датчика зон через один мультиплексор 4:1, адрес зоны --- два переключателя; демультиплексор показывает номер зоны лампой. \emph{Критерий приёмки:} при каждом из четырёх адресов выход повторяет датчик именно этой зоны, а изменение других зон не влияет; лампа номера горит только одна (вкладка Диаграммы).
  \item \stars{3}~\textbf{Проект: «Три источника».} Выберите одно из трёх 4-разрядных чисел $a,b,c$ по двухбитному коду $s$ (при $s=3$ --- $c$). \emph{Критерий приёмки:} при $s=0,1,2,3$ выход равен $a,b,c,c$ (вкладка Диаграммы, пример чисел любой).
  \item \stars{3}~\textbf{Проект: «Мажоритарный выбор».} Реализуйте на одном мультиплексоре $4:1$ функцию «2 из 3» $F=AB+AC+BC$. \emph{Критерий приёмки:} таблица выхода совпадает с таблицей истинности из главы~\ref{ch:comb} при всех 8 комбинациях.
\end{enumerate}
\end{taskbox}

\section{Что дальше}

Мы научились выбирать и распределять: из многих сигналов получать один и
обратно. Но выбирать можно не только датчики, но и \emph{результаты
вычислений}: пусть сначала посчитаются сумма, разность, логическое «И», а
мультиплексор по коду операции пропустит нужный ответ. Остался вопрос: как
схема из вентилей вообще \emph{считает}? В следующей главе
(глава~\ref{ch:arith}) мы построим сумматор, вычитатель, компаратор и соберём
из них, и из мультиплексора этой главы, арифметико-логическое устройство.
